\documentclass[11pt]{article}

\usepackage[T1]{fontenc}
\usepackage[utf8]{inputenc}
\usepackage[a4paper,margin=27mm]{geometry}
\usepackage{parskip}
\usepackage[hidelinks]{hyperref}

\pagestyle{empty}
\setlength{\parindent}{0pt}

\begin{document}

\begin{flushright}
\today
\end{flushright}

Editors of \emph{Array}\\
Elsevier

\textbf{Re: Regular Paper submission, ``Resource-Sensitive Certification of
Tool-Augmented Agent Swarms: Coalgebraic Semantics and LangGraph Validation''}

Dear Editors,

Please consider the enclosed manuscript as a Regular Paper for publication in
\emph{Array}. The paper studies how an execution trace from a tool-augmented
multi-agent workflow can be checked against explicit obligations for routing,
permissions, budgets, append-only evidence, provenance, audit freshness, and
support for final-answer claims. It combines a deterministic coalgebraic model
of observable swarm behaviour with a side-conditioned, resource-sensitive
certificate calculus whose labels are motivated by Subexponential Linear Logic
(SELL).

The contribution is deliberately bounded. Coalgebra describes the raw events
that a framework execution may emit; the certificate layer recognises those
finite paths for which the required ledger evidence is present. The paper proves
prefix closure and finite-run preservation results for the stated operational
fragment and gives an indexed, transactional checker for that fragment. It does
not claim a bialgebraic distributive law, full SELL proof search, mechanisation,
factual correctness of tool output, authenticated provenance, adversarial
security, or protocol-progress guarantees. Provenance is structural: the
checker establishes that a required record passes the documented
explicit-provenance predicate and is linked to an append-only write, not that
the record is externally true or cryptographically authenticated.

The accompanying evaluation separates four questions and reports the evidence
for each one:

\begin{itemize}
  \item a 19-trace conformance suite contains five accepted traces and fourteen
        targeted negative controls, all classified as expected;
  \item three deterministic workflows were executed with LangGraph 1.2.11 and
        their native version-2 streams were retained before conversion to the
        canonical trace schema; all three unmodified workflows were accepted;
  \item fourteen certification duties were mutated once in each of the three
        framework traces, and all 42 exact mutations were rejected at the
        mutated step with the expected violation code; and
  \item a scaling sanity check covers valid communication traces from 10 to
        10,000 events, using one warm-up and seven measured repetitions per
        size.
\end{itemize}

All evaluation data are original synthetic data generated for this study. The
workflows are deterministic, use no language model and no commercial API, and
exercise one orchestration framework. Accordingly, the experiments establish
executable conformance, mutation sensitivity, and integration feasibility in
the studied setting; they do not establish production-scale generality or
comparative superiority. The associated reproducibility artifact contains the
checker, LangGraph adapter and workflows, experiment runner, tests, pinned
dependencies, generated tables, environment metadata, and integrity manifests.

The manuscript fits \emph{Array} because it joins theoretical computer science
with executable software-engineering validation for data-intensive and
intelligent systems. Its formal contribution gives an explicit semantic and
resource-accounting boundary; its software contribution turns that boundary
into a reproducible validation pipeline over native framework records. The
article should therefore be relevant to readers working across formal methods,
software reliability, data provenance, workflow orchestration, and intelligent
systems infrastructure.

I confirm that the manuscript is original, has not been published, and is not
under consideration by another journal. As the sole author, I approve the
submission and take responsibility for the manuscript, software, data, and
reported results. I declare no competing interests. This research received no
specific grant from funding agencies in the public, commercial, or not-for-profit
sectors. The study used no human participants, personal data, animals, or
third-party research dataset.

During preparation of the work, I used OpenAI Codex to assist with repository
auditing, software development and testing, literature organisation, and
language and structural editing. I reviewed and edited all outputs,
independently verified the cited sources and reported results, and take full
responsibility for the content of the submission. The corresponding declaration
is included in the manuscript.

Thank you for considering this work for \emph{Array}.

Sincerely,

\vspace{2mm}
Carlos Ram\'irez Ovalle\\
Departamento de Ciencias Naturales y Matem\'aticas\\
Pontificia Universidad Javeriana Cali\\
Cali, Colombia\\
\href{mailto:carlosovalle@javerianacali.edu.co}{carlosovalle@javerianacali.edu.co}

\end{document}
